We chose the STM32L476RG Nucleo board as our benchmarking platform, featuring 96KB of RAM, 1MB of Flash memory,
a 32-bit ARM Cortex M4 which can be clocked up to 80MHz and a true random number generator.
Due to its large memory footprint and the limited amount of memory available on our target platform,
we take into consideration only the \texttt{hqc-128} instance.

We use \texttt{libopencm3} (1.0.0) as firmware library and compile the project with\newline \texttt{gcc-arm-none-eabi} (7.2.1)
using the \texttt{-O3} optimization flag. The board is clocked at a range of frequencies between 16 and 80 MHz, in order to
evaluate the impact of the memory wait states; the frequencies and corresponding amounts of wait cycles are reported in
the table below:
\begin{table}[ht]
    \centering
    \label{tab:waitstates}
    \begin{tabular}{|c|c|c|c|c|c|}
        \hline
        Frequency (MHz) & 16 & 32 & 48 & 64 & 80 \\
        \hline 
        Wait States & 0 & 1 & 2 & 3 & 4  \\
        \hline    
    \end{tabular}
\end{table}

\subsection{Execution time}
\pgfplotstableread[col sep=comma]{  % enc
    Label, Reference (R3), sparsemul (R3), Reference (R4), sparsemul (R4), FAFFT (R4)
    0,      2726,          18920,           136262,         18965,          8834
    1,      2748,          18937,           138227,         18991,          10107
    2,      2799,          18986,           142093,         19049,          11595
    3,      2886,          19064,           147971,         19154,          13189
    4,      3035,          19196,           155762,         19324,          14856
}\armEncTiming % enc

\pgfplotstableread[col sep=comma]{  % dec
    Label,               Reference (R3), sparsemul (R3), Reference (R4), sparsemul (R3), FAFFT (R4)
    0,                 1942,           9010,           68750,          8915,           5043
    1,                 1950,           9018,           69733,          8926,           5677
    2,                 1965,           9027,           71657,          8941,           6409
    3,                 1989,           9045,           74568,          8958,           7176
    4,                 2024,           9071,           78415,          8998,           7958
}\armDecTiming

\pgfplotstableread[col sep=comma]{  % encaps
    Label,           Reference (R3), sparsemul (R3), Reference (R4), sparsemul (R4), FAFFT (R4)
    0,                 3913,           20105,          137873,         20571,          10441
    1,                 3958,           20148,          139835,         20643,          11763
    2,                 4079,           20261,          143631,         20783,          13333
    3,                 4289,           20461,          149685,         21051,          15087
    4,                 4626,           20787,          157722,         21473,          16994
}\armEncapsTiming

\pgfplotstableread[col sep=comma]{  % decaps
    Label,               Reference (R3), sparsemul (R3), Reference (R4), sparsemul (R4), FAFFT (R4)
    0,                 5923,           29182,          206691,         29557,          15546
    1,                 5974,           29230,          209618,         29634,          17504
    2,                 6110,           29361,          215263,         29795,          19804
    3,                 6341,           29578,          224238,         30084,          22326
    4,                 6724,           29926,          236118,         30540,          25015
}\armDecapsTiming

\pgfplotstableread[col sep=comma]{  % mult
    Label,               Reference (R3), sparsemul (R3), Reference (R4), sparsemul (R4), FAFFT (R4)
    0,                 848,            7897,           66777,          7791,           3062
    1,                 848,            7897,           67747,          7791,           3688
    2,                 851,            7897,           69653,          7791,           4404
    3,                 856,            7897,           72537,          7791,           5146
    4,                 862,            7897,           76345,          7791,           5892
}\armMultTiming
  % contains data for the plots

To measure the execution times we use the Data Watchpoint and Trace Cycle Count Register (\texttt{DWT\_CYCCNT}),
a 32-bit register that counts the number of clock cycles elapsed since its last reset with a resolution
of one clock cycle. To avoid overflows, we reset the register after all the measurements of a single round
are completed; to reduce uncertainty, we average the measurements taken over 1000 rounds.
\subsection{Welch's t-test}
To evaluate the existence of an eventual side-channel vulnerability we use the Test Vector Leakage Assessment (TVLA)~\cite{schneider2015leakage},
based on Welch's t-test. The aim of the test is to provide a quantitative measure of how much the means of two
populations differ; the t-test statistic is computed as follows:
\begin{equation}
    t = \frac{\mu_0 - \mu_1}{\sqrt{\frac{s_0^2}{n_0} + \frac{s_1^2}{n_1}}}
\end{equation}  
where $\mu, s^2$ and $n$ are the mean, variance and number of samples of the two populations, respectively.
Small $t$ values give evidence to accept the null hypothesis that the two populations have the same mean; 
usually, for the sake of simplicity, a threshold of $|t| > 4.5$ is defined to reject the null hypothesis with
a confidence level of $99,999\%$.\\
In our case, we take the populations to be the execution times of an algorithm, where in the first we draw
new inputs randomly each time, while in the second we use the same inputs for all the rounds; each population
is composed of 1000 samples.\\